% ======================================================================
\chapter{Conclusão e próximos passos}
\label{cap:conclusao}
% ======================================================================

Em duas horas, o caminho foi do zero ao fluxo completo de um projeto no
GitHub: conta protegida e benefícios de estudante, Git e GitHub CLI
instalados, commits bem escritos, branches, conflitos resolvidos,
contribuição por fork e Pull Request, um site no ar pelo GitHub Pages, uma
automação no GitHub Actions, arquivos grandes com o Git LFS e commits
assinados.

Mais importante que decorar comandos é entender o modelo por trás deles: o
histórico é uma sequência de commits, uma branch é um ponteiro, e um Pull
Request é uma conversa antes de integrar. Com isso, os comandos novos passam a
fazer sentido sozinhos.

\section{Para continuar estudando}

\begin{itemize}
  \item \textbf{Pro Git}, de Scott Chacon e Ben Straub: o livro de referência
        do Git, gratuito e com tradução para o português,
        \url{https://git-scm.com/book/pt-br/v2}.
  \item \textbf{GitHub Skills}: cursos interativos dentro do próprio GitHub,
        \url{https://skills.github.com}.
  \item \textbf{Learn Git Branching}: exercícios visuais de branch e merge, em
        português, \url{https://learngitbranching.js.org/?locale=pt_BR}.
  \item \textbf{Documentação do GitHub}: \url{https://docs.github.com/pt}.
\end{itemize}

\begin{dica}[O próximo passo]
Escolha um projeto seu, mesmo pequeno, e coloque no GitHub com README,
licença e um workflow. Um repositório bem cuidado vale mais no portfólio do
que muitos repositórios vazios.
\end{dica}
